% Generated draft: requires mathematical and source review.
\documentclass[11pt]{article}
\usepackage[margin=25mm]{geometry}
\usepackage{fontspec}
\setmainfont{Latin Modern Roman}
\usepackage{amsmath,amssymb,mathtools,unicode-math}
\setmathfont{Latin Modern Math}
\usepackage{enumitem,needspace,xurl,hyperref}
\setlength{\parindent}{0pt}
\setlength{\parskip}{5pt}
\setlength{\emergencystretch}{3em}
\newcommand{\N}{\mathbb{N}}
\newcommand{\Z}{\mathbb{Z}}
\newcommand{\Q}{\mathbb{Q}}
\newcommand{\R}{\mathbb{R}}
\newcommand{\C}{\mathbb{C}}
\newcommand{\F}{\mathbb{F}}
\newcommand{\A}{\mathbb{A}}
\newcommand{\PP}{\mathbb{P}}
\newcommand{\E}{\mathbb{E}}
\newcommand{\eps}{\varepsilon}
\newcommand{\dd}{\,\mathrm{d}}
\begin{document}
\section*{UnsolvedMath ID: 30005427}
\textbf{Source problem code:} catalogue entry
\subsection{Definitions and mathematical statement}
The supplied excerpt introduces the additional assumption \(k=\C\), an algebra \(\Lambda\cong kQ/I\), an admissible ideal \(I\), and the number \(n\) of vertices of \(Q\). Here \(kQ\) denotes the path algebra of the quiver \(Q\). Under the usual convention, admissibility means that \(J^m\subseteq I\subseteq J^2\) for some integer \(m\geq2\), where \(J\) is the ideal generated by the arrows. The word ``moreover'' in the excerpt indicates earlier hypotheses that are not supplied.

For a finite-dimensional \(\Lambda\)-module \(M\), viewed as a representation of \(Q\), and a dimension vector \(\mathbf d=(d_1,\ldots,d_n)\in\Z_{\geq0}^n\), let \(\operatorname{Gr}_{\mathbf d}(M)\) denote the quiver Grassmannian parametrizing submodules \(N\subseteq M\) with \(\dim_k N_i=d_i\) at each vertex \(i\). Write \(\chi\) for the topological Euler characteristic and \(\mathbf y^{\mathbf d}=\prod_{i=1}^n y_i^{d_i}\). In this finite-dimensional setting, the displayed formula in the excerpt has the standard interpretation
\[
F(M)=\sum_{\mathbf d\in\Z_{\geq0}^n}\chi\!\left(\operatorname{Gr}_{\mathbf d}(M)\right)\mathbf y^{\mathbf d}\in\Z[y_1,\ldots,y_n].
\]
The excerpt says ``any \(\Lambda\)-module'' without specifying finite dimensionality; the definition above explains the usual finite-dimensional interpretation and does not resolve the intended module domain.

\textbf{Reported conjecture, with unresolved formulation.} The original excerpt asserts that a map \(\Phi\) is conjectured to be an isomorphism. The supplied clean statement describes it as a parametrization map associated with the \(F\)-polynomials above and asks whether
\[
\Phi\text{ is an isomorphism}.
\]
However, the record gives neither the definition of \(\Phi\), its source and target, nor the structures with respect to which it should be an isomorphism. Moreover, the conjecture sentence precedes the heading ``4. Parametrization'' and the assumption \(k=\C\); the excerpt does not establish that this later paragraph supplies the hypotheses or construction for that conjecture. Consequently, a precise quantified conjecture cannot be recovered from this record.
\subsection{Short English statement}
The record asks whether an unspecified map \(\Phi\) is an isomorphism and associates it with module \(F\)-polynomials. It supplies an \(F\)-polynomial formula over \(\C\), but omits the map, its domain, its codomain, and potentially earlier hypotheses. This is an incomplete statement-only draft requiring human mathematical review.
\subsection{Sources}
\begin{itemize}
\item Local UnsolvedMath registry, supplied source record with permanent ID \(30005427\), title ``Module Parametrization by \(F\)-Polynomials'', exact source problem code \texttt{catalogue entry}. The record attributes the problem to a 2023 Oberwolfach report; no individual proposer is supplied.
\item \emph{Representation Theory of Quivers and Finite-Dimensional Algebras} (2023), Oberwolfach Reports, pp.~397--486, DOI \texttt{10.4171/owr/2023/7}. Supplied source URL and literature-source link labeled ``original OWR report'': \url{https://doi.org/10.4171/owr/2023/7}.
\item Supplied literature assessment, checked \textbf{2026-08-22}: the record labels the problem open but states that \(\Phi\), its source, and its target are undefined in the extracted record, so no resolution can be verified from that extraction. This is supplied dated curation, not an independently verified literature assessment; its time-sensitive status and cited claims require independent checking.
\item The registry also supplies the statement-status label \texttt{corrected\_verified} and a claim that the clean statement was checked against the official report. Those are source-provided verification claims, not verification performed for this draft, and they do not supply the missing mathematical data.
\end{itemize}
\end{document}
